\section{Positive Hodge diagonals and the Kodaira--Nadel route}

The positive-twist argument uses the decomposition of Deligne--Illusie
\cite{DI87}.  For a smooth projective variety of dimension $n$, put
\[
 P_t(m):\quad H^{t-i}(X,\Omega_X^i\otimes L^{\otimes m})=0
 \quad\hbox{for every }0\leq i\leq n,\qquad t>n.
\]
Only nonnegative cohomology degrees occur.  Each degree $t-i$ is
positive, which is essential for the application of Serre vanishing.
No finite-dimensionality assertion is needed.

\begin{lemma}[Finite type is preserved by tensoring with a line]
\label{lem:kn-line-tensor-finite-type}
\lean{EnvelopeProperty.SchemeInvertibleTensorFiniteType.isFiniteType_tensorRight}
\leanok
Let $M$ be a quasicoherent finite-type module on a scheme $X$, and let
$L$ be an invertible module.  Then $M\otimes L$ is of finite type.
\end{lemma}
\begin{proof}
Refine a trivializing cover of $L$ by affine opens.  On each such open
the restricted module $M$ has a finite generating family.
Restriction of the tensor product, the chosen line frame and the
tensor unit comparison identify the restricted $M\otimes L$ with
the restricted $M$.  Transport its finite generating family through
this isomorphism.  These are finite local generators on an actual
affine open cover of $X$.
\end{proof}

\begin{theorem}[Serre vanishing with an arbitrary coherent fixed factor]
\label{thm:kn-coherent-ample-serre}
\lean{EnvelopeProperty.SchemeAmpleCoherentSerreVanishing.exists_sheafCohomology_bound}
\uses{lem:kn-line-tensor-finite-type,thm:fc-closed-projective-serre-vanishing}
\leanok
Let $X$ carry a geometrically ample invertible module $L$ over a
noetherian ring, and let $M$ be a quasicoherent finite-type module.
There is an integer $N$ such that
$H^q(X,L^{\otimes m}\otimes M)=0$ for every $m\geq N$ and every $q>0$.
\end{theorem}
\begin{proof}
Choose $r>0$ and a closed projective embedding such that the pullback
of $\mathcal O(1)$ is $L^{\otimes r}$.  For each $0\leq a<r$, the
module $M\otimes L^{\otimes a}$ is quasicoherent and of finite type.
Closed-projective Serre vanishing gives a threshold $N_a$ for that
module, simultaneously in every positive cohomology degree.  Put
$T=\max_a N_a$.  If $m\geq rT$, write $m=a+rd$, where $0\leq a<r$
and $d\geq T$.  The tensor-power and embedding comparisons identify
the corresponding Serre twist with $L^{\otimes m}\otimes M$.
Transport vanishing through this module isomorphism.
\end{proof}

\begin{corollary}[One bound for a finite family]
\label{cor:kn-finite-family-serre}
\lean{EnvelopeProperty.SchemeAmpleCoherentSerreVanishing.exists_uniform_sheafCohomology_bound}
\uses{thm:kn-coherent-ample-serre}
\leanok
For a finite family of quasicoherent finite-type modules $M_i$, one
threshold $N$ annihilates every $H^q(X,L^{\otimes m}\otimes M_i)$
for every $i$, every $m\geq N$, and every $q>0$.
\end{corollary}
\begin{proof}
Apply the theorem to each member and take the maximum of the finitely
many resulting thresholds.  Every exponent beyond this maximum is
beyond the threshold for each member.  For an empty family take zero.
\end{proof}

\begin{lemma}[Finite type of the actual exterior cotangent sheaves]
\label{lem:kn-exterior-finite-type}
\lean{EnvelopeProperty.Vanishing.SmoothExteriorFiniteType.isFiniteType}
\leanok
Let $X$ be smooth of relative dimension $n$ over a field.  Every
exterior cotangent sheaf $\Omega_X^i$, including $i=0$ and $i>n$,
is of finite type.
\end{lemma}
\begin{proof}
Around each point choose an affine standard-smooth chart.  Its
cotangent module has a basis indexed by $\{1,\ldots,n\}$.  In exterior
degree $i$, the wedges indexed by the $i$-element subsets form a
finite basis.  The affine sheaf comparison and restriction along the
chart isomorphism transport this basis to a finite free frame on the
original open site.  Its generator sections are finite local
generators.  If $i>n$, the indexing set is empty, which gives the zero
module and is still finite.
\end{proof}

\begin{corollary}[Eventual vanishing on all upper Hodge diagonals]
\label{cor:kn-hodge-serre}
\lean{EnvelopeProperty.Vanishing.HodgeSerreVanishing.exists_bound,EnvelopeProperty.Vanishing.HodgeSerreVanishing.exists_upper_diagonal_bound}
\uses{cor:kn-finite-family-serre,lem:kn-exterior-finite-type}
\leanok
Let $X$ be smooth of dimension $n$ over a field, with an ample line
$L$.  There is one $N$ such that $P_t(m)$ holds for every $m\geq N$
and every $t>n$.
\end{corollary}
\begin{proof}
The finitely many modules $\Omega_X^0,\ldots,\Omega_X^n$ are
quasicoherent and of finite type.  Apply the finite-family Serre bound.
For $i\leq n<t$, the degree $t-i$ is positive, so each of the groups
in $P_t(m)$ vanishes.  The same bound works for every $t>n$ because
the Serre bound already works in every positive cohomology degree.
\end{proof}

\begin{lemma}[Amplification of an eventual property]
\label{lem:kn-frobenius-diagonal-iteration}
\lean{EnvelopeProperty.Vanishing.FrobeniusIteration.of_eventually}
\leanok
Let $p\geq2$ and let $P(m)$ be a property of nonnegative integers
such that $P(pm)$ implies $P(m)$ for every $m>0$.  If $P(m)$ holds
for every sufficiently large $m$, then it holds for every $m>0$.
\end{lemma}
\begin{proof}
Induction shows that $P(mp^r)$ implies $P(m)$ for every $r$.  Since $p\geq2$,
induction gives $r+1\leq p^r$.  Choose $r$ at least an eventual
threshold.  Since $m\geq1$, $mp^r$ is beyond that threshold.
Thus $P(mp^r)$ holds, and the iterated implication gives $P(m)$.
\end{proof}

\begin{corollary}[Conditional upper-diagonal descent criterion]
\label{cor:kn-upper-diagonal-descent}
\lean{EnvelopeProperty.Vanishing.HodgeSerreVanishing.upper_diagonal_of_descent}
\uses{cor:kn-hodge-serre,lem:kn-frobenius-diagonal-iteration}
\leanok
In the preceding geometric setting, fix $p\geq2$ and $t>n$.
Assume that $P_t(pm)$ implies $P_t(m)$ for every $m>0$.
Then $P_t(m)$ holds for every $m>0$.
\end{corollary}
\begin{proof}
The actual Hodge-Serre theorem supplies eventual $P_t(m)$.
Apply the amplification lemma with precisely the assumed implication.
\end{proof}

\begin{remark}[Remaining geometric comparisons]
\label{rmk:kn-geometric-frontier}
The iteration lemma is independent of Deligne--Illusie.  To
produce the one-step implication, the geometric proof must identify the cohomology
of the Frobenius twist with the original cohomology, apply the
projection formula to the derived decomposition, and prove the
vanishing of bounded hypercohomology from a zero entire diagonal.
The decomposition itself is temporarily treated as an input from
Deligne--Illusie.  These comparisons are separate obligations.
Passing to characteristic zero further requires spreading out and
cohomological specialization.  The local and Kawamata--Viehweg
vanishing arguments and the multiplier-ideal comparison are needed
before this route proves Nadel vanishing; ordinary eventual Serre
vanishing alone does not prove it.
\end{remark}
